\chapter{Tecnologie e strumenti}
\label{chap:tecnologie}


\section{Linguaggi di programmazione}
\textbf{PHP} è un linguaggio di programmazione general purpose particolarmente diffuso nello sviluppo di applicazioni web lato server. È stato scelto, in versione 8.3, come linguaggio 
principale per il backend di \textit{ErpAI} grazie alla sua maturità, alla vasta disponibilità di librerie e strumenti per lo sviluppo web e alla sua integrazione con il framework 
Laravel. Un ulteriore vantaggio è la possibilità di sviluppare applicazioni strutturate e facilmente manutenibili attraverso le funzionalità offerte dalle versioni moderne 
del linguaggio.


\section{Framework e librerie}
\subsection{Laravel}
\textbf{Laravel} è un framework PHP open source progettato per lo sviluppo di applicazioni web. Fornisce una struttura organizzata per la realizzazione del backend e 
mette a disposizione funzionalità per la gestione delle richieste HTTP, del routing, dell'accesso al database, dell'autenticazione e delle operazioni asincrone. 
È stato scelto, in versione 11.x, per la sua ampia dotazione di funzionalità integrate, per la facilità di sviluppo e manutenzione del codice e per l'integrazione con gli altri componenti 
utilizzati nel progetto.

\subsection{Filament}
\textbf{Filament} è un framework per Laravel dedicato alla realizzazione di pannelli amministrativi e interfacce gestionali. Permette di costruire rapidamente risorse, tabelle, 
form e pagine attraverso componenti già pronti e fortemente integrati con Laravel. È stato scelto, in versione 3.x, per semplificare la realizzazione dell'interfaccia gestionale di \textit{ErpAI}, 
riducendo il codice necessario per implementare le funzionalità di amministrazione e mantenendo una stretta integrazione con il backend.

\subsection{Livewire}
\textbf{Livewire} è un framework full-stack per Laravel che consente di creare interfacce web dinamiche utilizzando prevalentemente codice PHP. Il suo principale 
vantaggio è la possibilità di aggiornare dinamicamente il DOM comunicando con il server senza dover gestire manualmente uno stato complesso lato client via JavaScript. 
Essendo la tecnologia su cui si basa Filament, la sua adozione garantisce la massima coerenza nello sviluppo dei componenti interattivi del gestionale.

\section{Tecnologie di persistenza e gestione dei dati}
\subsection{MariaDB}
\textbf{MariaDB} è un sistema di gestione di database relazionale open source, compatibile con SQL. È stato scelto, in versione 10.11, come database dell'applicazione per la sua affidabilità, 
le buone prestazioni e l'integrazione con Laravel tramite il sistema di astrazione del database fornito dal framework. Il modello relazionale risulta inoltre particolarmente 
adatto alla gestione delle numerose entità e relazioni presenti in un gestionale aziendale.

\subsection{Redis}
\textbf{Redis} è un sistema di gestione dati in memoria utilizzabile come database, cache e sistema di gestione delle code. La sua caratteristica principale è la possibilità di 
effettuare operazioni molto rapide grazie alla gestione dei dati prevalentemente in memoria. Nel progetto viene utilizzato, nella versione 7, per la gestione della cache e delle code, consentendo 
di separare le operazioni più onerose dal normale ciclo di elaborazione delle richieste.

\subsection{Laravel Horizon}
\textbf{Laravel Horizon} è un componente dell'ecosistema Laravel dedicato al monitoraggio e alla gestione delle code basate su Redis. Fornisce un'interfaccia per controllare i 
processi in esecuzione, le attività completate e quelle fallite, oltre a informazioni sulle prestazioni delle code. È stato scelto, nella versione 5.x, per semplificare la gestione delle operazioni 
asincrone del progetto e permettere di monitorarne lo stato direttamente dall'applicazione.

\section{Ambiente di sviluppo e containerizzazione}

\subsection{Docker latest}
\textbf{Docker} è una piattaforma di containerizzazione che permette di eseguire applicazioni e servizi all'interno di ambienti isolati chiamati container. I container consentono 
di distribuire insieme all'applicazione le dipendenze necessarie, riducendo le differenze tra gli ambienti di sviluppo ed esecuzione. Nel progetto viene utilizzato come base 
dell'ambiente di sviluppo e per l'esecuzione isolata dei diversi servizi.

\subsection{DDEV latest}
\textbf{DDEV} è un ambiente di sviluppo locale, basato su Docker, progettato per semplificare la configurazione di applicazioni web e dei relativi servizi. Utilizza la 
containerizzazione per fornire agli sviluppatori un ambiente riproducibile e isolato. È stato scelto per rendere uniforme l'ambiente di sviluppo tra i membri del team e 
semplificare la configurazione dei servizi necessari all'esecuzione di \textit{ErpAI}.

\section{Strumenti di qualità e testing}
\subsection{Pest}
\textbf{Pest} è un framework per il testing automatizzato in ambiente PHP. Progettato con un approccio incentrato sulla sintassi fluida e l'espressività del codice, 
permette di scrivere unit e integration test estremamente leggibili e sintetici rispetto al classico PHPUnit. È stato scelto per verificare automaticamente il corretto 
funzionamento delle funzionalità del backend di \textit{ErpAI} e prevenire l'insorgere di regressioni durante l'evoluzione del codice.

\subsection{PHPStan e Larastan}
\textbf{PHPStan} è uno strumento di analisi statica per PHP che permette di individuare potenziali errori e incongruenze nel codice senza doverlo eseguire. L'analisi viene 
effettuata considerando principalmente i tipi e le possibili combinazioni dei dati utilizzati dal programma. \textbf{Larastan} estende PHPStan fornendo regole e funzionalità 
specifiche per l'ecosistema Laravel, consentendo di analizzare correttamente costrutti e convenzioni tipiche del framework. Sono stati scelti per aumentare l'affidabilità e 
la qualità del codice, individuando preventivamente problematiche che potrebbero manifestarsi durante l'esecuzione dell'applicazione.

